\documentclass[10pt,a4paper]{amsart}
\usepackage{amsmath,amssymb,amsthm,graphicx}
% Retained minimal preamble for the Stacks Project chapter sites-cohomology.tex (Cohomology on Sites).
% Verbatim excerpt of the chapter preamble: Xy-pic setup (native preamble lines 14-18) and the native
% theorem declarations with the mathematical macros (native preamble lines 154-217). Copied mechanically;
% nothing added, removed or redefined. The wrapper class and the amsmath/amssymb/amsthm/graphicx base
% packages are supplied by the assembler, so the class/externaldocument part of the native preamble is omitted.
% For commutative diagrams we use Xy-pic
\usepackage[all]{xy}

% We use 2cell for 2-commutative diagrams.
\xyoption{2cell}
% Theorem environments.
%
\theoremstyle{plain}
\newtheorem{theorem}[subsection]{Theorem}
\newtheorem{proposition}[subsection]{Proposition}
\newtheorem{lemma}[subsection]{Lemma}

\theoremstyle{definition}
\newtheorem{definition}[subsection]{Definition}
\newtheorem{example}[subsection]{Example}
\newtheorem{exercise}[subsection]{Exercise}
\newtheorem{situation}[subsection]{Situation}

\theoremstyle{remark}
\newtheorem{remark}[subsection]{Remark}
\newtheorem{remarks}[subsection]{Remarks}

\numberwithin{equation}{subsection}

% Macros
%
\def\lim{\mathop{\mathrm{lim}}\nolimits}
\def\colim{\mathop{\mathrm{colim}}\nolimits}
\def\Spec{\mathop{\mathrm{Spec}}}
\def\Hom{\mathop{\mathrm{Hom}}\nolimits}
\def\Ext{\mathop{\mathrm{Ext}}\nolimits}
\def\SheafHom{\mathop{\mathcal{H}\!\mathit{om}}\nolimits}
\def\SheafExt{\mathop{\mathcal{E}\!\mathit{xt}}\nolimits}
\def\Sch{\mathit{Sch}}
\def\Mor{\mathop{\mathrm{Mor}}\nolimits}
\def\Ob{\mathop{\mathrm{Ob}}\nolimits}
\def\Sh{\mathop{\mathit{Sh}}\nolimits}
\def\NL{\mathop{N\!L}\nolimits}
\def\CH{\mathop{\mathrm{CH}}\nolimits}
\def\proetale{{pro\text{-}\acute{e}tale}}
\def\etale{{\acute{e}tale}}
\def\QCoh{\mathit{QCoh}}
\def\Ker{\mathop{\mathrm{Ker}}}
\def\Im{\mathop{\mathrm{Im}}}
\def\Coker{\mathop{\mathrm{Coker}}}
\def\Coim{\mathop{\mathrm{Coim}}}

% Boxtimes
%
\DeclareMathSymbol{\boxtimes}{\mathbin}{AMSa}{"02}

%
% Macros for moduli stacks/spaces
%
\def\QCohstack{\mathcal{QC}\!\mathit{oh}}
\def\Cohstack{\mathcal{C}\!\mathit{oh}}
\def\Spacesstack{\mathcal{S}\!\mathit{paces}}
\def\Quotfunctor{\mathrm{Quot}}
\def\Hilbfunctor{\mathrm{Hilb}}
\def\Curvesstack{\mathcal{C}\!\mathit{urves}}
\def\Polarizedstack{\mathcal{P}\!\mathit{olarized}}
\def\Complexesstack{\mathcal{C}\!\mathit{omplexes}}
% \Pic is the operator that assigns to X its picard group, usage \Pic(X)
% \Picardstack_{X/B} denotes the Picard stack of X over B
% \Picardfunctor_{X/B} denotes the Picard functor of X over B
\def\Pic{\mathop{\mathrm{Pic}}\nolimits}
\def\Picardstack{\mathcal{P}\!\mathit{ic}}
\def\Picardfunctor{\mathrm{Pic}}
\def\Deformationcategory{\mathcal{D}\!\mathit{ef}}
\usepackage{cleveref}
\crefname{originallemma}{original lemma}{original lemmas}
\Crefname{originallemma}{Original lemma}{Original lemmas}
\makeatletter
\newlabel{lemma-K-flat-resolution}{{06Y4}{source}{}{sourcepointer-1}{}}
\newlabel{lemma-K-flat-resolution@cref}{{[originallemma][0][]06Y4}{[1][0][]source}{}{}{}}
\makeatother
\crefname{originallemma}{original lemma}{original lemmas}
\Crefname{originallemma}{Original lemma}{Original lemmas}
\makeatletter
\newlabel{lemma-K-flat-two-out-of-three-ses}{{0BYH}{source}{}{sourcepointer-2}{}}
\newlabel{lemma-K-flat-two-out-of-three-ses@cref}{{[originallemma][0][]0BYH}{[1][0][]source}{}{}{}}
\makeatother
\crefname{originallemma}{original lemma}{original lemmas}
\Crefname{originallemma}{Original lemma}{Original lemmas}
\makeatletter
\newlabel{lemma-bounded-flat-K-flat}{{06YD}{source}{}{sourcepointer-3}{}}
\newlabel{lemma-bounded-flat-K-flat@cref}{{[originallemma][0][]06YD}{[1][0][]source}{}{}{}}
\makeatother
\crefname{originallemma}{original lemma}{original lemmas}
\Crefname{originallemma}{Original lemma}{Original lemmas}
\makeatletter
\newlabel{lemma-last-one-flat}{{0653}{source}{}{sourcepointer-4}{}}
\newlabel{lemma-last-one-flat@cref}{{[originallemma][0][]0653}{[1][0][]source}{}{}{}}
\makeatother
\crefname{originallemma}{original lemma}{original lemmas}
\Crefname{originallemma}{Original lemma}{Original lemmas}
\makeatletter
\newlabel{lemma-pullback-K-flat}{{06YC}{source}{}{sourcepointer-5}{}}
\newlabel{lemma-pullback-K-flat@cref}{{[originallemma][0][]06YC}{[1][0][]source}{}{}{}}
\makeatother
\title{Tor-amplitude bounded below is equivalent to being representable by a K-flat complex of flat modules vanishing below the bound}
\author{The Stacks Project Authors}
\date{stacks-project:sites-cohomology.tex:local-snapshot-20261011; GNU FDL 1.2 or later}
\begin{document}
\maketitle
\noindent\textit{Editorial note.} All delivered mathematical text is the author's own, copied from the pinned Stacks Project chapter file sites-cohomology.tex; no statement or proof step was rewritten, reordered or removed.\par\medskip
\noindent\textit{Editorial note.} The counted proof is the author's complete argument: the proof is written out in full: the choice of the K-flat resolution, the flatness of the truncation cokernel via the last-one-flat criterion, the K-flatness of the kernel and the truncated complex by two-out-of-three, and the pullback argument giving the uniformity stated in the last sentence.\par\medskip
\noindent\textit{Editorial note.} The label lemma-bounded-below-tor-amplitude names the counted statement, and the proof environment immediately following it in the source is delivered verbatim as its proof.\par\medskip
\noindent\textit{Editorial note.} The statement and proof are copied from the pinned chapter file; the five cited inputs are kept as external pointers with their canonical Stacks tags.\par\medskip
\noindent\textit{Editorial note.} The preamble is the chapter's own retained minimal preamble (sources/stacks-sites-cohomology/retained-preamble.tex), a verbatim excerpt of the native preamble: the Xy-pic setup and the native theorem declarations with the mathematical macros.\par\medskip
\noindent\textit{Allowed standard prerequisites (editorial).}\begin{itemize}
\item K-flat complexes of modules on a ringed site and the two-out-of-three property. Same chapter, the notion of a K-flat complex, the existence of K-flat resolutions with flat terms, and the two-out-of-three lemma for short exact sequences of complexes. Supplies the representative K\textasciicircum{}bullet chosen at the start and the two applications of two-out-of-three showing L\textasciicircum{}bullet and E\textasciicircum{}bullet are K-flat.
\item Flatness criterion on the last term of a bounded above complex. Same chapter, the lemma concluding flatness of the cokernel of a map of flat modules from the vanishing of negative tensor cohomology. Supplies flatness of N = Coker(K\textasciicircum{}\{a-1\} -> K\textasciicircum{}a) and of the image I in K\textasciicircum{}a.
\item Pullback of a K-flat complex is K-flat. Same chapter, the lemma showing that the pullback of a K-flat complex of flat modules along a morphism of ringed sites is again K-flat with flat terms. Supplies the final assertion by pulling the whole construction back along an arbitrary morphism of ringed sites.
\end{itemize}\medskip
\ifdefined\hypertarget\hypertarget{sourcepointer-1}{}\fi
\noindent\textit{Source reference.} Original lemma 06Y4: sources/stacks-sites-cohomology/pinned-original-source.tex, source lines 3207--3207. Prerequisite pointer invoked at proof line(s) 12489; it names a standard fact of the same chapter that the counted argument uses, and every step of the counted proof built on it is retained verbatim in the delivered body.\par\medskip
\ifdefined\hypertarget\hypertarget{sourcepointer-2}{}\fi
\noindent\textit{Source reference.} Original lemma 0BYH: sources/stacks-sites-cohomology/pinned-original-source.tex, source lines 3097--3097. Prerequisite pointer invoked at proof line(s) 12521; it names a standard fact of the same chapter that the counted argument uses, and every step of the counted proof built on it is retained verbatim in the delivered body.\par\medskip
\ifdefined\hypertarget\hypertarget{sourcepointer-3}{}\fi
\noindent\textit{Source reference.} Original lemma 06YD: sources/stacks-sites-cohomology/pinned-original-source.tex, source lines 3120--3120. Prerequisite pointer invoked at proof line(s) 12519; it names a standard fact of the same chapter that the counted argument uses, and every step of the counted proof built on it is retained verbatim in the delivered body.\par\medskip
\ifdefined\hypertarget\hypertarget{sourcepointer-4}{}\fi
\noindent\textit{Source reference.} Original lemma 0653: sources/stacks-sites-cohomology/pinned-original-source.tex, source lines 12403--12403. Prerequisite pointer invoked at proof line(s) 12498; it names a standard fact of the same chapter that the counted argument uses, and every step of the counted proof built on it is retained verbatim in the delivered body.\par\medskip
\ifdefined\hypertarget\hypertarget{sourcepointer-5}{}\fi
\noindent\textit{Source reference.} Original lemma 06YC: sources/stacks-sites-cohomology/pinned-original-source.tex, source lines 3508--3508. Prerequisite pointer invoked at proof line(s) 12526; it names a standard fact of the same chapter that the counted argument uses, and every step of the counted proof built on it is retained verbatim in the delivered body.\par\medskip
\begin{lemma}
\label{lemma-bounded-below-tor-amplitude}
Let $(\mathcal{C}, \mathcal{O})$ be a ringed site. Let $E$ be an object of
$D(\mathcal{O})$. Let $a \in \mathbf{Z}$. The following
are equivalent
\begin{enumerate}
\item $E$ has tor-amplitude in $[a, \infty]$.
\item $E$ can be represented by a K-flat complex $\mathcal{E}^\bullet$
of flat $\mathcal{O}$-modules with $\mathcal{E}^i = 0$ for
$i \not \in [a, \infty]$.
\end{enumerate}
Moreover, we can choose $\mathcal{E}^\bullet$ such that any pullback
by a morphism of ringed sites is a K-flat complex with flat terms.
\end{lemma}

\begin{proof}
The implication (2) $\Rightarrow$ (1) is immediate. Assume (1) holds.
First we choose a K-flat complex $\mathcal{K}^\bullet$
with flat terms representing $E$, see Lemma \ref{lemma-K-flat-resolution}.
For any $\mathcal{O}$-module $\mathcal{M}$ the cohomology of
$$
\mathcal{K}^{n - 1} \otimes_\mathcal{O} \mathcal{M} \to
\mathcal{K}^n \otimes_\mathcal{O} \mathcal{M} \to
\mathcal{K}^{n + 1} \otimes_\mathcal{O} \mathcal{M}
$$
computes $H^n(E \otimes_\mathcal{O}^\mathbf{L} \mathcal{M})$.
This is always zero for $n < a$. Hence if we apply
Lemma \ref{lemma-last-one-flat} to the complex
$\ldots \to \mathcal{K}^{a - 1} \to \mathcal{K}^a \to \mathcal{K}^{a + 1}$
we conclude that $\mathcal{N} = \Coker(\mathcal{K}^{a - 1} \to \mathcal{K}^a)$
is a flat $\mathcal{O}$-module. We set
$$
\mathcal{E}^\bullet = \tau_{\geq a}\mathcal{K}^\bullet =
(\ldots \to 0 \to \mathcal{N} \to \mathcal{K}^{a + 1} \to \ldots )
$$
The kernel $\mathcal{L}^\bullet$ of
$\mathcal{K}^\bullet \to \mathcal{E}^\bullet$ is the complex
$$
\mathcal{L}^\bullet = (\ldots \to \mathcal{K}^{a - 1} \to
\mathcal{I} \to 0 \to \ldots)
$$
where $\mathcal{I} \subset \mathcal{K}^a$ is the image of
$\mathcal{K}^{a - 1} \to \mathcal{K}^a$.
Since we have the short exact sequence
$0 \to \mathcal{I} \to \mathcal{K}^a \to \mathcal{N} \to 0$
we see that $\mathcal{I}$ is a flat $\mathcal{O}$-module.
Thus $\mathcal{L}^\bullet$ is a bounded
above complex of flat modules, hence K-flat by
Lemma \ref{lemma-bounded-flat-K-flat}.
It follows that $\mathcal{E}^\bullet$ is K-flat by
Lemma \ref{lemma-K-flat-two-out-of-three-ses}.

\medskip\noindent
Proof of the final assertion. Let
$f : (\mathcal{C}', \mathcal{O}') \to (\mathcal{C}, \mathcal{O})$
be a morphism of ringed sites. By Lemma \ref{lemma-pullback-K-flat}
the complex $f^*\mathcal{K}^\bullet$ is K-flat with flat terms.
The complex $f^*\mathcal{L}^\bullet$ is K-flat as it is a bounded
above complex of flat $\mathcal{O}'$-modules. We have a short exact
sequence of complexes of $\mathcal{O}'$-modules
$$
0 \to f^*\mathcal{L}^\bullet \to f^*\mathcal{K}^\bullet \to
f^*\mathcal{E}^\bullet \to 0
$$
because the short exact sequence
$0 \to \mathcal{I} \to \mathcal{K}^a \to \mathcal{N} \to 0$
of flat modules pulls back to a short exact sequence.
By Lemma \ref{lemma-K-flat-two-out-of-three-ses}.
the complex $f^*\mathcal{E}^\bullet$ is K-flat and the proof is complete.
\end{proof}

\end{document}
